\begingroup
\begin{tikzpicture}[
  font=\thesisfigurefont,
  box/.style={draw=black!55,rounded corners=2pt,align=center,inner sep=4pt,line width=.6pt},
  wide/.style={box,text width=12.7cm,minimum height=1.15cm,fill=blue!3},
  flow/.style={-{Latex[length=2mm]},draw=black!65,line width=.65pt},
  node distance=6mm
]
\node[wide,fill=black!3] (train) {\textbf{联邦持续分类训练}\par 客户端本地更新与子空间构建\quad 服务器模型聚合与子空间融合\par 每个类别阶段结束后得到全局模型};
\node[wide,below=of train] (submit) {\textbf{提交模型或预测}\par 选择任务顺序、完成阶段及测试分类头选择方式\par 阶段全局模型与分类预测按任务和样本标识对应};
\node[wide,below=of submit] (test) {\textbf{MedCL：任务与阶段评价}\par 同一阶段全局模型\quad 按测试任务选择分类头\quad 以准确率构建阶段--任务矩阵};
\node[wide,below=of test,fill=green!4] (view) {\textbf{任务性能、持续学习能力与资源}\par ACC\quad BWTR\quad RMA\quad MPE\quad DRR\par 客户端分布与参与率\quad 模型和子空间通信开销};
\draw[flow] (train) -- (submit);
\draw[flow] (submit) -- (test);
\draw[flow] (test) -- (view);
\end{tikzpicture}
\endgroup
